%!TeX root = ../main.tex
% ---------------------------------------------------------------% ---------------------------------------------------------------
\section{Discussion}
\label{sec:discussion}
% ---------------------------------------------------------------

While recursive models incorporate potentially useful information about reference quality, there is no definitive recursive approach to ranking a set of units. Uncertainties and errors in metadata, combined with flexibility in ranking parameters, lead to many different potential rankings for the same units and reference lists: some are perhaps more reliable than others. Consequently, comparisons of our spectral rankings with other rankings are not only a face-validity check, but also an exploration of the range of variation due to methodology. Section~\ref{sec:indicators} compared $v$ with established lists and indicators. Here we look more closely at expert judgment, in the 2025 Journal Quality List of the Australian Business Deans Council \textcite[][ABDC]{australianbusinessdeanscouncil2025ABDCJournal2025}. It lists 219 A* journals, of which 135 are included in the OAS corpus. The difference in coverage reflects our registry filter that removed journals having substantial content in fields such as psychology, sociology, law, education, technology and public policy. Of the 135 journals that overlap, 98 lie among the 135 highest-$v$ corpus journals, and all but five lie among corpus journals with $v>1$. The five low-$v$ journals cover business and economic history, hospitality and tourism, information technology and occupational psychology, fields where experts' views could understandably allocate high importance to a journal.

Ritzberger \citep{ritzbergerRankingJournalsEconomics2008} presented a spectral ranking of a selection of economics, business, industrial relations and statistics journals from the 2006 version of the Clarivate Journal Citation Reports. He classified 25 journals as A+ or A. Of these, one is not in our corpus, 23 have $v>1$ in the baseline model and one (in mathematical economics) has $v<1$ with strong in-and-out coupling to journals outside our corpus. Three of the A+/A journals, the finance journals \textit{Journal of Finance}, \textit{Review of Financial Studies} and \textit{Journal of Financial Economics}, are business journals in our classification, and the remainder are economics. The paucity of business journals in this list has been discussed by \textcite{albersMisleadingRankingsResearch2009} (see also comments by \textcite{ritzbergerMisleadingRankingsResearch2009}). 

While several other drivers for the $E$/$B$ separation have been discussed, the comparison between Figure~\ref{fig:combined_vs_separate} and Figure~\ref{fig:eb_alignment} implicates the clique of top $E$ journals. \textcite{ritzbergerRankingJournalsEconomics2008} anticipated our counterfactual by considering the ranking with his top five journals omitted (these are also our top five), finding like us that the accounting and marketing journals rise in cardinal rank while maintaining their relative positions. He concluded that `rankings within minor subfields will largely agree with the restrictions from the overall ranking – a methodological advantage of weight-based ranking methods'. We are uncertain about this assessment of the effect of the clique. Ritzberger's JCR journal list from 2002 had 261 journals, while our 2024 JCR has 1433: the selection criteria are slightly different, but the main change is the addition of many items from the Emerging Journals Citation Index around 2023. The additional journals combined with variable publication growth rates over the intervening two decades prevents direct comparisons between the two sets of rankings. However, the comparison between the high-ranked journals in our Figure~\ref{fig:combined_vs_separate} and Figure~\ref{fig:eb_alignment} implies that either (a) artificial removal of top $E$ journals or (b) recirculation and fan-out through institutions radically intersperses $E$ and $B$ journals and thus changes ordinal ranks in the combined sets. This occurs despite order preservation within the sets of $E$ and $B$ journals. It is clear that claims about intrinsic inter-field merit derived from combined rankings can rest on structural features of the citation network -- implying that discussion might move to consider the merit of influential cliques. \textcite{biehlRelationshipsAcademicBusiness2006}, \textcite{fourcadeSuperiorityEconomists2015}, \textcite{gibsonWhichJournalRankings2014}, \textcite{aistleitnerCitationPatternsEconomics2019} and \textcite{bacciniExploringShapeEconomics2026} present more detailed treatments of the internal and coupling structure of $E$/$B$ subfields consistent with this conclusion.  The sensitivity of overall spectral rankings to the presence of subfields and other structural differences has been emphasised by \textcite{pinskiCitationInfluenceJournal1976}, \textcite{waltmanRelationEigenfactorAudience2010}, \textcite{heilFieldDependentNaturePageRank2023} and \textcite{cavallaroSensitivityCentralityMetrics2024}.

The critique of unit spectral ranking by \textcite{vaccarioWhenStandardNetwork2022} exposes the risk of over-interpreting a unit citation matrix in terms of the transfer of intellectual influence. On the other hand, the distinction between popularity and prestige at the unit level \citep{bollenJournalStatus2006} remains valid and useful. Our preliminary study of enclaves suggests that combining influence based on unit-level spectral ranking with citation prominence based on work-level reference analysis can reveal interesting features of the science system.  If it were possible to recursively rank the sparse directed acyclic graph (DAG) of work-work references without introducing the arbitrary parameters found in the Katz-Hubbell/PageRank or HRank \citep{caoRankingAcademicInstitutions2023} methods, it might be feasible to construct a cardinal ranking of works carrying ideas through actual reference trails \citep[see also][]{golosovskyCitationAnalysisDynamics2019}.

\textcite{franceschetTenGoodReasons2010} advances ten reasons to use Eigenfactor metrics, providing a convenient structure for a critical summary of our work. For example, the reasons claim there is value in adjusting for reference intensity, the duration of the target window, and the presence of journal self-citations. We show these measures change rank scores by $\approx$ 0.2 (s.d. in the log), and rank values by $\approx$ 5\%. A decision to make these adjustments would rest on the purpose of the ranking, but all could be seen as advantageous. Other reasons, such as a solid mathematical background, an intuitive stochastic interpretation, a sound axiomatic foundation, and relatability to other ranking metrics are clearly favorable. Perhaps less clearly advantageous is a single ranking over the entire citation network. We have seen that relative rankings cannot be compared reliably without paying close attention to the structure of the citation network, in particular the capacity of small, tight cliques to harbour influence probability. Large scale rankings such as Clarivate's Eigenfactor and Article Influence Metrics should be rescaled when used to compare communities with different publication and referencing characteristics \citep{yuLogarithmicEigenfactorSolving2021}. The use of reference weights to recursively establish journal influence is beneficial in some applications, but the influence values require careful interpretation due to `fictitious' links arising from projection as well as structural drivers in the network. As Franceschet notes, spectral ranking  is readily implemented and extended, and combined with a large scale, open and free metadata resource such as OpenAlex it opens the way to low-cost, auditable rankings with minimal reliance on commercial interests. The cost is small: from the edge list, the baseline ranking takes about 2.6~s on one processor core, against 0.7~s for citation counts and 1.2~s for MNCS on the same corpus (Supplement~1, Table~\ref{S1-tab:cost}), and a bootstrap replicate takes about the same time as the baseline, so 1,000 replicates run in under an hour. However, omitted from Franceschet's list are the response of spectral ranking to errors and uncertainties. Our analysis suggests that for small metadata error levels, bibliometric spectral rankings are amenable to confidence intervals, while large errors remain attached to the affected unit and do not necessarily propagate widely.

\subsection{Limitations and conditions of use}
\label{sec:limitations}

The bipartite fixed point has a simple reading. Researchers at influential institutions reference works in influential journals, and influential journals publish works that are referenced from influential institutions. The eigenvector is the fixed point of this mutual recognition, the bibliometric counterpart of the hubs and authorities of \textcite{kleinberg1999authoritative}. Within Social Systems Citation Theory it is one of the mechanisms by which the science system stabilises itself, and from the side of the individual researcher it resembles the Matthew effect \citep{mertonMatthewEffectScience1968}. The ranking shows the fixed point; it does not demonstrate either mechanism. Who the anchors of recognition are is defined by the eigenvector itself, which is why the comparisons of Section~\ref{sec:indicators} with rankings built in other ways matter.

This reading also indicates where the model is useful. It suits fields in which journals and institutions both carry evaluative weight and a journal list sets the corpus boundary, comparisons within a field or across neighbouring fields with the clique structure reported, and the descriptive mapping of influence flows. It does not suit the ranking of whole universities across many fields, because an institution node carries only the institution's works in the corpus. It is fragile for small units: because $v=A\pi/a_p$, a small work count amplifies $v$, and the highest-ranked institution in the baseline, the Federal Reserve Bank of Minneapolis, has an author-fractional work count of 44 over five years. It should not be used as a single number without its interval, nor for fields whose literature lies largely outside journals.

The model has structural costs as well. Institution nodes pool fields; the split-node test (Section~\ref{sec:fields}) shows that the effect is small here, a few hundredths in $\log v$, but it need not be small in other fields, and splitting nodes by field is a direct remedy. Some nodes are affiliations rather than employers: the National Bureau of Economic Research and the Centre for Economic Policy Research receive a fraction of many works through second affiliations and rank second and fourth among institutions; such bodies can be reported separately. Combined rankings depend on dense cliques (Section~\ref{sec:fields}) and on the corpus boundary: including the sentinel units changes $\log v$ by 0.3 for journals (s.d.). Identity errors in the metadata concentrate on institutions and have to be corrected unit by unit (Supplement~2). Finally, the self-reinforcing loop that the eigenvector describes is itself a reason for caution, since rankings used for evaluation feed back into the behaviour they measure \citep{espelandRankingsReactivityHow2007}. We therefore present $v$, with its error band, as a measure of recognition within the reference network rather than of quality.

